\documentclass[11pt]{article}
\usepackage[margin=1in]{geometry}
\usepackage{booktabs}
\usepackage{graphicx}
\usepackage{hyperref}
\usepackage{enumitem}
\setlist{nosep}
\graphicspath{{../plots/out/}}

\title{MBLA: Task-Scoped Permission Inference for BoundClaw\\\large Overview of benchmark findings}
\author{BoundClaw MBLA working group}
\date{\today}

\begin{document}
\maketitle

\section{What MBLA is}
MBLA is the policy-inference and recovery-admission component of BoundClaw. At each delegation it proposes the
capability and confinement permissions the receiver needs, choosing only from what the delegator already holds.
The trusted checks (attenuation on the delegator, ceiling on the receiver) bound every proposal, so a wrong or
manipulated model can cost utility but never safety. When a receiver is denied mid-task it may request one
concrete permission; MBLA decides whether to grant it.

The Go package (\texttt{mbla/}) is what the gVisor fork imports. The benchmark (\texttt{bench/}) reproduces every
number here. Design decisions, assumptions and findings are logged in \texttt{research/}, including a section-by-section
mapping onto the paper (\texttt{research/paper\_alignment.md}).

\section{Setup}
\begin{itemize}
  \item \textbf{Tasks.} The paper's three workflows (BC-1 GitHub triage, BC-2 build and staging deploy, BC-3
    workspace-to-Slack): 17 task templates, 288 test items and 120 dev items, with injection variants. The test
    split was frozen by hash before any test run.
  \item \textbf{References.} Two people labeled every candidate permission independently (Table~\ref{tab:audit});
    35 disagreements were resolved by rule. Agreement is reported on the independent labels.
  \item \textbf{Methods.} Decision models (Qwen3-Reranker 0.6B/4B/8B, Laya 0.4B, CLM-v0.1-8B) answer
    ``is this permission needed?'' per candidate. Generative LLMs (Qwen3 0.6B--32B, Qwen3-30B-A3B, Phi-4-mini,
    gpt-oss-20b/120b) output the permission set as JSON. Baselines: full parent, operator manifest, keyword match.
  \item \textbf{Protocol.} Thresholds tuned on dev only with a pre-registered objective; 95\% confidence intervals
    from a cluster bootstrap over test templates; paired tests with Holm correction. Latency measured separately,
    each model alone on an L40S GPU (E5).
\end{itemize}

\begin{table}[h]\centering\small\begin{tabular}{lr}
\toprule
Reference audit & Value \\
\midrule
Rows (17 templates) & 185 \\
Agreement, Alvi vs Joseph & 0.941 \\
Cohen's kappa (95\% CI) & 0.853 [0.749, 0.941] \\
Krippendorff's alpha (humans) & 0.853 \\
Rows needing resolution & 35 \\
\bottomrule
\end{tabular}

\caption{Reference audit: independent agreement between two human reviewers.}\label{tab:audit}\end{table}

\section{Findings}

\paragraph{1. The paper's recovery rule lets every attack through; scored admission fixes it.}
Under the paper's rule (grant any single action inside the recovery bound), a malicious receiver reaches every
attack the parent allows: 100\% for all methods (Figure~\ref{fig:recovery}). Requiring the request to also be in the
scorer's \emph{admit} set cuts this to 13\% for Qwen3-8B (95\% CI 8--20\%) while 97\% of legitimate tasks still
succeed (attack reduction significant, $p < 0.001$; success change not significant). gpt-oss and Reranker-8B land
near 50\%. Methods with coarse 0/1 scores gain nothing, because their dev tuning admits everything.

\begin{figure}[h]\centering\includegraphics[width=\linewidth]{recovery.pdf}
\caption{Share of reachable attacks after recovery (cap 3). Error bars: 95\% CI.}\label{fig:recovery}\end{figure}

\paragraph{2. Mid-size generative models are the most precise; decision models are faster but over-grant.}
Qwen3-32B, Qwen3-30B, Qwen3-8B and Qwen3-4B keep excess capability authority below 0.1 while missing under 7\% of
needed permissions (Table~\ref{tab:quality}). gpt-oss-120b is not more precise than Qwen3-8B. Decision models answer
in tens of milliseconds but, zero-shot, grant 0.38--0.85 of the avoidable authority (Figure~\ref{fig:quality}).

\begin{figure}[h]\centering\includegraphics[width=0.85\linewidth]{quality_latency.pdf}
\caption{Excess capability authority (E1, test) against median latency (E5, exclusive GPU).}\label{fig:quality}\end{figure}

\begin{table}[h]\centering\small\begin{tabular}{lrrrrr}
\toprule
Method & Type & p50 ms & EAC (C) & MAC (C) & Exact \\
\midrule
manifest-only & baseline & -- & 0.025 & 0.185 & 0.000 \\
qwen3-32b-awq & generative & 829 & 0.061 & 0.038 & 0.389 \\
qwen3-30b & generative & 280 & 0.073 & 0.066 & 0.131 \\
qwen3-8b & generative & 777 & 0.086 & 0.031 & 0.330 \\
qwen3-4b & generative & 436 & 0.092 & 0.036 & 0.233 \\
gpt-oss-20b & generative & 4112 & 0.135 & 0.036 & 0.330 \\
gpt-oss-120b & generative & 2538 & 0.139 & 0.022 & 0.366 \\
qwen3-reranker-8b & decision & 187 & 0.379 & 0.058 & 0.024 \\
qwen3-reranker-0.6b & decision & 28 & 0.397 & 0.072 & 0.007 \\
lexical & baseline & -- & 0.473 & 0.000 & 0.000 \\
qwen3-1.7b & generative & 231 & 0.636 & 0.000 & 0.000 \\
qwen3-reranker-4b & decision & 118 & 0.642 & 0.016 & 0.066 \\
clm-zs & decision & 91 & 0.709 & 0.028 & 0.000 \\
qwen3-0.6b & generative & 131 & 0.745 & 0.010 & 0.000 \\
phi-4-mini & generative & 382 & 0.794 & 0.024 & 0.000 \\
laya & decision & 26 & 0.848 & 0.088 & 0.000 \\
set-only & bound & -- & 1.000 & 0.000 & 0.000 \\
\bottomrule
\end{tabular}

\caption{E1 test results. EAC and MAC: excess and missing capability authority (paper \S6.3). Latency from E5.}\label{tab:quality}\end{table}

\paragraph{3. External validation on FORTIS.}
Scored with FORTIS's own metric code, Reranker-8B matches 30B--120B generators on skill selection (exact match
0.45--0.47), while tool selection is hard for every method (best 0.29). Decision-model thresholds tuned on our
data do not transfer: they either abstain or grant everything (Figure~\ref{fig:fortis}, Table~\ref{tab:fortis}).

\begin{figure}[h]\centering\includegraphics[width=\linewidth]{fortis.pdf}
\caption{FORTIS exact match per method (task 1: 579 queries, task 2: 1,543 queries).}\label{fig:fortis}\end{figure}

\begin{table}[h]\centering\small\begin{tabular}{lrrrr}
\toprule
Model & Right skill & Riskier skill & Right tools & Extra tools \\
\midrule
gpt-oss-120b & 47\% & 35\% & 29\% & 33\% \\
qwen3-30b & 45\% & 36\% & 18\% & 45\% \\
qwen3-reranker-8b & 45\% & 39\% & 11\% & 19\% \\
gpt-oss-20b & 45\% & 38\% & 25\% & 39\% \\
qwen3-4b & 43\% & 41\% & 16\% & 45\% \\
qwen3-32b-awq & 41\% & 49\% & 21\% & 48\% \\
qwen3-reranker-4b & 40\% & 48\% & 9\% & 46\% \\
qwen3-8b & 36\% & 49\% & 17\% & 46\% \\
capmas-scores & 35\% & 52\% & 0\% & 100\% \\
bge-large-zs & 34\% & 51\% & 0\% & 100\% \\
qwen3-reranker-0.6b & 32\% & 58\% & 6\% & 74\% \\
qwen3-1.7b & 25\% & 43\% & 0\% & 99\% \\
clm-zs & 18\% & 47\% & 0\% & 100\% \\
laya & 7\% & 85\% & 0\% & 100\% \\
capmas & 3\% & 94\% & 0\% & 99\% \\
phi-4-mini & 2\% & 93\% & 0\% & 100\% \\
set-only & 2\% & 96\% & 0\% & 100\% \\
qwen3-0.6b & 1\% & 96\% & 0\% & 100\% \\
lexical & 0\% & 0\% & 0\% & 100\% \\
\bottomrule
\end{tabular}

\caption{FORTIS results with the benchmark's official classification.}\label{tab:fortis}\end{table}

\paragraph{4. What the scorer sees, and how the backbone is used.}
For Qwen3-8B, adding the operator manifest to the input cuts excess authority from about 0.23 to 0.09
(Figure~\ref{fig:inputs}). On the same Qwen3-8B weights, generating the set (EAC 0.09) is far more precise than a
reranker head (0.38) or CLM's contrastive heads (0.71). Operating curves show generative scores are coarse, while
decision models trade excess for misses smoothly (Figure~\ref{fig:curves}).

\begin{figure}[h]\centering
\includegraphics[width=0.49\linewidth]{inputs.pdf}\hfill\includegraphics[width=0.49\linewidth]{curves.pdf}
\caption{Left: input ablation (E3). Right: operating curves (E2).}\label{fig:inputs}\label{fig:curves}\end{figure}

\paragraph{5. Correctness of the trusted checks.}
The Go attenuation checker agrees with an independent set-semantics implementation on 14,480 generated
parent/child pairs and chains up to 8 hops (0 false accepts, 0 false rejects); a deliberately broken checker is caught.

\section{Status}
\begin{itemize}
  \item Done: library, references, E1--E7, E9, FORTIS, significance tests, LLM judge labels (an independent
    judge's adjustments change excess by at most 0.03 and leave the ranking unchanged).
  \item Pending: optional human validation of the judge;
    optional fine-tuned CLM (E1b).
\end{itemize}

\end{document}
